% Part: sets-functions-relations
% Chapter: relations
% Section: relations-as-sets
% Hindi translation (hi-Deva-IN) of Open Logic commit 9620cc73f9c8e0ad003c514a5d3748f29611c4c0.

\documentclass[../../../include/open-logic-section]{subfiles}

\begin{document}

\olfileid[hi]{sfr}{rel}{set}
\olsection{समुच्चयों के रूप में संबंध}

\begin{explain}
\olref[sfr][set][imp]{sec} में हमने कुछ महत्त्वपूर्ण समुच्चयों का उल्लेख
किया था: $\Nat$, $\Int$, $\Rat$, $\Real$। इनमें से कुछ समुच्चयों के
!!{element}ों के बीच के कई रोचक संबंध निस्संदेह आपको याद होंगे। उदाहरण के
लिए, इनमें से प्रत्येक समुच्चय पर एक सर्वथा मानक \emph{क्रम-संबंध} है। इसके
अतिरिक्त \emph{तत्सम होने} का संबंध भी है, जिसमें प्रत्येक वस्तु स्वयं से
संबंधित है और किसी अन्य वस्तु से नहीं। आगे हमें और भी कई रोचक संबंध मिलेंगे,
और संभावित संबंधों की संख्या तो इससे भी अधिक है। फिर भी, उनका पुनरवलोकन
करने से पहले हम यह स्पष्ट करने से आरम्भ करेंगे कि संबंधों को एक विशेष प्रकार
के समुच्चय के रूप में देखा जा सकता है।
  
इसके लिए \olref[sfr][set][pai]{sec} की दो बातें याद कीजिए। पहली, किसी
\emph{क्रमित युग्म} की संकल्पना याद कीजिए: $a$ और $b$ दिए होने पर हम
~$\tuple{a, b}$ बना सकते हैं। महत्त्वपूर्ण बात यह है कि यहाँ अवयवों का क्रम
\emph{वास्तव में} मायने रखता है। अतः यदि $a \neq b$, तो $\tuple{a, b} \neq \tuple{b, a}$।
(इसकी तुलना अक्रमित युग्मों, अर्थात् $2$-अवयवी
समुच्चयों, से कीजिए, जहाँ $\{a, b\}=\{b, a\}$।) दूसरी, \emph{कार्तीय
गुणनफल} की संकल्पना याद कीजिए: यदि $A$ और $B$ समुच्चय हैं, तो हम
~$A \times B$ बना सकते हैं, जो उन सभी युग्मों $\tuple{x, y}$ का समुच्चय है
जिनमें $x \in A$ और $y \in B$। विशेष रूप से, $A^{2}= A \times A$,~$A$
से लिए गए सभी क्रमित युग्मों का समुच्चय है।
  
अब हम किसी समुच्चय पर एक विशेष संबंध पर विचार करेंगे: $<$-संबंध, जो
प्राकृत संख्याओं के समुच्चय~$\Nat$ पर है। संख्याओं के उन सभी युग्मों $\tuple{n, m}$
के समुच्चय पर विचार कीजिए जिनमें $n<m$, अर्थात्,
\[
  R=\Setabs{\tuple{n, m}}{n, m \in \Nat \text{ और } n<m}.
\]
$n$ का $m$ से छोटा होना और युग्म $\tuple{n, m}$ का $R$ का अवयव होना आपस
में निकटता से जुड़े हैं; ठीक-ठीक:
\[
      n<m\text{ तभी और केवल तभी }\tuple{n, m} \in R.
\]
वास्तव में, किसी भी सूचना को खोए बिना हम समुच्चय $R$ को $<$-संबंध
\emph{ही} मान सकते हैं, जो $\Nat$ पर है।

इसी प्रकार संख्याओं के बीच किसी भी संबंध के लिए हम $\Nat^{2}$ का एक
उपसमुच्चय बना सकते हैं। इसके विपरीत, संख्या-युग्मों का कोई भी समुच्चय
$S \subseteq \Nat^{2}$ दिया हो, तो संख्याओं के बीच उसके अनुरूप एक संबंध
होता है: $n$ का $m$ से वह संबंध तभी और केवल तभी होता है जब $\tuple{n, m} \in S$।
इससे निम्न परिभाषा उचित ठहरती है:
\end{explain}
  
\begin{defn}[द्वयी संबंध]
समुच्चय $A$ पर एक \emph{द्वयी संबंध}~$A^{2}$ का कोई उपसमुच्चय है। यदि $R
\subseteq A^{2}$,~$A$ पर एक द्वयी संबंध है और $x, y \in A$, तो हम
कभी-कभी $Rxy$ (अथवा $xRy$) लिखते हैं, जिसका अर्थ $\tuple{x, y} \in R$ है।
\end{defn}
  
\begin{ex}
  \ollabel{relations}
प्राकृत संख्याओं के युग्मों के समुच्चय $\Nat^{2}$ को इस प्रकार एक
द्विविमीय सारणी में सूचीबद्ध किया जा सकता है:
\[
  \begin{array}{ccccc}
  \mathbf{\tuple{ 0,0 }} & \tuple{ 0,1 } &
    \tuple{ 0,2 } & \tuple{ 0,3 } & \ldots\\
  \tuple{ 1,0 } & \mathbf{\tuple{ 1,1 }} &
    \tuple{ 1,2 } & \tuple{ 1,3 } & \ldots\\
  \tuple{ 2,0 } & \tuple{ 2,1 } &
    \mathbf{\tuple{ 2,2 }} & \tuple{ 2,3 } & \ldots\\
  \tuple{ 3,0 } & \tuple{ 3,1 } & \tuple{ 3,2 } &
    \mathbf{\tuple{ 3,3 }} & \ldots\\
  \vdots & \vdots & \vdots & \vdots & \mathbf{\ddots}
  \end{array}
\]
यहाँ हमने विकर्ण को मोटे अक्षरों में रखा है, क्योंकि विकर्ण पर स्थित युग्मों
से बना $\Nat^2$ का उपसमुच्चय, अर्थात्,
\[
  \{\tuple{0,0 }, \tuple{ 1,1 }, \tuple{ 2,2 }, \dots\},
  \]
~$\Nat$ पर \emph{तत्समक संबंध} है। (चूँकि तत्समक संबंध का प्रायः उपयोग
होता है, हम $\Id{A}=\Setabs{\tuple{ x,x }}{x \in
A}$ को किसी भी समुच्चय $A$ के लिए परिभाषित करते हैं।) विकर्ण के ऊपर स्थित सभी युग्मों का
उपसमुच्चय, अर्थात्,
\[
  L = \{\tuple{ 0,1 },\tuple{ 0,2 },\ldots,\tuple{ 1,2 },
  \tuple{ 1,3 }, \dots, \tuple{ 2,3 }, \tuple{ 2,4 },\ldots\},
\]
\emph{‘से छोटा है’ संबंध} है, अर्थात् $Lnm$ तभी और केवल तभी जब $n<m$।
विकर्ण के नीचे स्थित युग्मों का उपसमुच्चय, अर्थात्,
\[
  G=\{ \tuple{ 1,0 },\tuple{ 2,0 },\tuple{
    2,1 }, \tuple{ 3,0 },\tuple{ 3,1 },\tuple{ 3,2 }, \dots\},
\]
\emph{‘से बड़ा है’ संबंध} है, अर्थात् $Gnm$ तभी और केवल तभी जब $n>m$।
$L$ और $I$ का सम्मिलन, जिसे हम $K=L\cup I$ कह सकते हैं, \emph{‘से छोटा
या बराबर है’ संबंध} है: $Knm$ तभी और केवल तभी जब $n \le m$। इसी प्रकार,
$H=G \cup
I$, \emph{‘से बड़ा या बराबर है’ संबंध} है। ये संबंध $L$, $G$,
$K$, और $H$ विशेष प्रकार के संबंध हैं, जिन्हें \emph{क्रम} कहा जाता है।
$L$ और $G$ का यह गुणधर्म है कि कोई भी संख्या स्वयं से $L$ या $G$ संबंध
नहीं रखती (अर्थात् प्रत्येक $n$ के लिए न तो $Lnn$ और न ही $Gnn$)। इस
गुणधर्म वाले संबंधों को \emph{अस्वतुल्य} कहते हैं; यदि वे क्रम भी हों, तो
उन्हें \emph{प्रबल क्रम} कहते हैं।
\end{ex}

\begin{explain}
यद्यपि क्रम और तत्समक संबंध महत्त्वपूर्ण तथा स्वाभाविक संबंध हैं, इस बात
पर बल देना चाहिए कि हमारी परिभाषा के अनुसार $A^{2}$ का \emph{कोई भी}
उपसमुच्चय~$A$ पर एक संबंध है, चाहे वह कितना भी अस्वाभाविक या कृत्रिम क्यों
न लगे। विशेष रूप से, $\emptyset$ किसी भी समुच्चय पर एक संबंध है
(\emph{रिक्त संबंध}, जिसमें कोई भी अवयव-युग्म संबंधित नहीं होता), और
$A^{2}$ स्वयं भी~$A$ पर एक संबंध है (जिसमें प्रत्येक युग्म संबंधित होता
है); इसे \emph{सार्वत्रिक संबंध} कहते हैं। परंतु
$E=\Setabs{\tuple{n, m}}{n>5 \text{ या } m \times n \ge 34}$ जैसा कोई
समुच्चय भी संबंध माना जाता है।
\end{explain}
  
\begin{prob}
  संबंध $\subseteq$ के !!{element}ों की सूची बनाइए; यह संबंध समुच्चय
  $\Pow{\{a, b, c\}}$ पर है।
\end{prob}

\end{document}
